\chapter{Hypothesis 1: Subwavelength Translation in Lateral, Vertical, and Diagonal Directions}
\label{ch:hyp1}

\Cref{sec:meth-resolution} showed that two \emph{stationary} scatterers
cannot be resolved as separate amplitude peaks below roughly half a
wavelength of separation. This chapter asks the closely related time-lapse
question: if a \emph{single} scatterer moves between a baseline and a
monitor survey by a sub-wavelength amount, is that movement visible in the
\emph{difference} of the two migrated amplitude images, and --- if not ---
can the 2D phase-plane estimator of \cref{sec:meth-phaseplane} recover it
instead? The displacement is swept from $2\lambda$ down to
$\tfrac{1}{32}\lambda$, using the same three migration algorithms throughout,
and the question is tested independently in three directions: lateral
(\cref{sec:hyp1-lateral}), vertical (\cref{sec:hyp1-vertical}), and diagonal
(\cref{sec:hyp1-diagonal}). \Cref{sec:hyp1-phaseplane} then applies the
global 2D weighted least-squares (WLS) phase-plane fit to all three
directions, and \cref{sec:hyp1-h15} explores an optional, local alternative
to that global fit.

\paragraph{Hypothesis 1.} Multi-dimensional phase-plane regression allows
tracking the subwavelength translation of a scatterer in the lateral,
vertical, and diagonal directions, using both amplitude- and phase-based
inference.

\paragraph{Hypothesis 1.5 (Optional).} Tracking the local phase gradients
$\partial\varphi/\partial x$ and $\partial\varphi/\partial y$ gives an
equivalent, simpler alternative to the global 2D phase-plane fit for
inferring horizontal and vertical translation.

\section{Lateral Movement}
\label{sec:hyp1-lateral}

A single PEC cylinder ($r=\SI{28}{mm}$, depth \SI{0.676}{m}) is held fixed
in the baseline survey and displaced laterally by each of eight scenarios in
the monitor survey, from $2\lambda$ down to $\tfrac{1}{32}\lambda$
(\cref{fig:tl-setup}), using the same gprMax domain and grid as
\cref{sec:meth-resolution}.

\begin{figure}[htbp]
  \centering
  \begin{subfigure}[t]{0.48\textwidth}
    \includegraphics[width=\linewidth]{TL_001_TimeLapse_Study__Model_Geometry__domain_4010_m_Δx__1_mm_PML.png}
    \caption{Model geometry}
  \end{subfigure}
  \hfill
  \begin{subfigure}[t]{0.48\textwidth}
    \includegraphics[width=\linewidth]{TL_002_Moving_scatterer_s1__all_8_scenarios__r__28_mm_depth__0676_m.png}
    \caption{All 8 displacement scenarios}
  \end{subfigure}
  \caption{Forward-model setup for the lateral time-lapse study: (a) the
    gprMax domain and grid; (b) the lateral displacement of scatterer
    \texttt{s1} across all eight scenarios, $r=\SI{28}{mm}$, depth
    \SI{0.676}{m}.}
  \label{fig:tl-setup}
\end{figure}

\Cref{fig:tl-bscans} shows the background-baseline and time-lapsed raw
B-scans, the background-subtracted result, and (for later use in
\cref{ch:hyp3}) the effect of adding synthetic Laplace-distributed noise at
$10\,\%$ of the signal standard deviation.

\begin{figure}[htbp]
  \centering
  \begin{subfigure}[t]{0.31\textwidth}
    \includegraphics[width=\linewidth]{TL_003_GPR_B-Scans__Background_Baseline_and_TimeLapsed_Models.png}
    \caption{Background \& time-lapsed}
  \end{subfigure}
  \hfill
  \begin{subfigure}[t]{0.31\textwidth}
    \includegraphics[width=\linewidth]{TL_004_GPR_B-Scans__Background_Subtracted.png}
    \caption{Background-subtracted}
  \end{subfigure}
  \hfill
  \begin{subfigure}[t]{0.31\textwidth}
    \includegraphics[width=\linewidth]{TL_005_GPR_B-Scans__With_Synthetic_Laplace_Noise_10_of_signal_std.png}
    \caption{With Laplace noise}
  \end{subfigure}
  \caption{Raw B-scans for the lateral time-lapse study: (a) background and
    time-lapsed models; (b) background-subtracted; (c) with synthetic
    Laplace-distributed noise at $10\,\%$ of the signal standard deviation
    (used in \cref{ch:hyp3}).}
  \label{fig:tl-bscans}
\end{figure}

The same tapering and $t_0$-shift conditioning as \cref{sec:meth-resolution}
(\cref{sec:meth-conditioning}) is applied before migration
(\cref{fig:tl-taper}).

\begin{figure}[htbp]
  \centering
  \begin{subfigure}[t]{0.48\textwidth}
    \includegraphics[width=\linewidth]{TL_006_Effect_of_Tapering_and_t0_Shift__2λ_dataset_single_trace.png}
    \caption{Single trace}
  \end{subfigure}
  \hfill
  \begin{subfigure}[t]{0.48\textwidth}
    \includegraphics[width=\linewidth]{TL_007_B-scan_effect_of_tapering_and_t0_shift__2λ_dataset.png}
    \caption{Full B-scan}
  \end{subfigure}
  \caption{Effect of tapering and the $t_0$ shift on the $2\lambda$ lateral
    displacement dataset: (a) a single trace; (b) the full B-scan.}
  \label{fig:tl-taper}
\end{figure}

\Cref{fig:tl-kirchhoff,fig:tl-gazdag,fig:tl-backprop} show, respectively,
the Kirchhoff-, Gazdag-, and back-propagation-migrated images for all
displacement scenarios, and the time-lapse difference (migrated monitor
minus migrated baseline) that isolates the moved scatterer.

\begin{figure}[htbp]
  \centering
  \begin{subfigure}[t]{0.48\textwidth}
    \includegraphics[width=\linewidth]{TL_008_Kirchhoff_Migration__All_8_Datasets____f_c15_GHz____aperture.png}
    \caption{All 8 scenarios}
  \end{subfigure}
  \hfill
  \begin{subfigure}[t]{0.48\textwidth}
    \includegraphics[width=\linewidth]{TL_009_Kirchhoff_Migration_zoomed____f_c15_GHz____aperture40.png}
    \caption{Zoomed}
  \end{subfigure}

  \vspace{0.6em}
  \begin{subfigure}[t]{0.48\textwidth}
    \includegraphics[width=\linewidth]{TL_010_Kirchhoff_Migration__TimeLapse_Differences_migrated__migrate.png}
    \caption{Time-lapse difference}
  \end{subfigure}
  \hfill
  \begin{subfigure}[t]{0.48\textwidth}
    \includegraphics[width=\linewidth]{TL_011_Kirchhoff_Migration__TimeLapse_Differences_zoomed.png}
    \caption{Difference, zoomed}
  \end{subfigure}
  \caption{Kirchhoff migration of the lateral time-lapse study
    ($f_c=\SI{15}{\GHz}$, aperture $=40$): (a,b) migrated image for all
    scenarios and zoomed; (c,d) migrated-monitor-minus-migrated-baseline
    difference and zoomed.}
  \label{fig:tl-kirchhoff}
\end{figure}

\begin{figure}[htbp]
  \centering
  \begin{subfigure}[t]{0.48\textwidth}
    \includegraphics[width=\linewidth]{TL_012_Gazdag_Phase-Shift_Migration__All_7_Datasets____f_c15_GHz.png}
    \caption{All scenarios}
  \end{subfigure}
  \hfill
  \begin{subfigure}[t]{0.48\textwidth}
    \includegraphics[width=\linewidth]{TL_013_Gazdag_Phase-Shift_Migration_zoomed____f_c15_GHz.png}
    \caption{Zoomed}
  \end{subfigure}

  \vspace{0.6em}
  \begin{subfigure}[t]{0.48\textwidth}
    \includegraphics[width=\linewidth]{TL_014_Gazdag_Migration__TimeLapse_Differences_migrated__migrated_b.png}
    \caption{Time-lapse difference}
  \end{subfigure}
  \hfill
  \begin{subfigure}[t]{0.48\textwidth}
    \includegraphics[width=\linewidth]{TL_015_Gazdag_Migration__TimeLapse_Differences_zoomed.png}
    \caption{Difference, zoomed}
  \end{subfigure}
  \caption{Gazdag phase-shift migration of the lateral time-lapse study
    ($f_c=\SI{15}{\GHz}$): (a,b) migrated image and zoomed; (c,d) time-lapse
    difference and zoomed.}
  \label{fig:tl-gazdag}
\end{figure}

\begin{figure}[htbp]
  \centering
  \begin{subfigure}[t]{0.48\textwidth}
    \includegraphics[width=\linewidth]{TL_016_Back-Propagation_E__All_7_Datasets____focus_at_1906_ns.png}
    \caption{$\lVert\mathbf{E}\rVert$, all scenarios}
  \end{subfigure}
  \hfill
  \begin{subfigure}[t]{0.48\textwidth}
    \includegraphics[width=\linewidth]{TL_017_Back-Propagation_E_zoomed____focus_at_1906_ns.png}
    \caption{$\lVert\mathbf{E}\rVert$, zoomed}
  \end{subfigure}

  \vspace{0.6em}
  \begin{subfigure}[t]{0.48\textwidth}
    \includegraphics[width=\linewidth]{TL_018_Back-Propagation_Ez__All_7_Datasets____focus_at_1906_ns.png}
    \caption{$E_z$, all scenarios}
  \end{subfigure}
  \hfill
  \begin{subfigure}[t]{0.48\textwidth}
    \includegraphics[width=\linewidth]{TL_019_Back-Propagation_Ez_zoomed____focus_at_1906_ns.png}
    \caption{$E_z$, zoomed}
  \end{subfigure}

  \vspace{0.6em}
  \begin{subfigure}[t]{0.48\textwidth}
    \includegraphics[width=\linewidth]{TL_020_Back-Propagation__TimeLapse_Differences_Ez_Ez__Ez_baseline.png}
    \caption{$E_z$ time-lapse difference}
  \end{subfigure}
  \hfill
  \begin{subfigure}[t]{0.48\textwidth}
    \includegraphics[width=\linewidth]{TL_021_Back-Propagation__TimeLapse_Differences_Ez_zoomed.png}
    \caption{Difference, zoomed}
  \end{subfigure}
  \caption{Time-reversal back-propagation migration of the lateral
    time-lapse study, focused at $t=\SI{1906}{\ns}$: field-magnitude (a,b)
    and $E_z$ (c,d) images, and the $E_z$ time-lapse difference (e,f).}
  \label{fig:tl-backprop}
\end{figure}

\subsection*{Amplitude-Based Detectability Summary}
\label{sec:tl-detectability}

\Cref{fig:tl-summary} overlays the signed time-lapse-difference amplitude
from all three algorithms and plots the normalised lateral point-spread
function of that difference at the true scatterer depth, as a function of
displacement.

\begin{figure}[htbp]
  \centering
  \begin{subfigure}[t]{0.48\textwidth}
    \includegraphics[width=\linewidth]{TL_022_TimeLapse_Migration_Comparison__Signed_Amplitude____f_c15_GH.png}
    \caption{Signed-amplitude comparison}
  \end{subfigure}
  \hfill
  \begin{subfigure}[t]{0.48\textwidth}
    \includegraphics[width=\linewidth]{TL_023_Normalised_Lateral_PSF__TimeLapse_Difference_at_True_Scatter.png}
    \caption{Normalised lateral PSF of the difference}
  \end{subfigure}
  \caption{(a) Signed time-lapse-difference amplitude for all three
    migration algorithms at $f_c=\SI{15}{\GHz}$; (b) the normalised lateral
    PSF of the difference image at the true scatterer depth, swept across
    lateral displacements from $2\lambda$ to $\tfrac{1}{32}\lambda$.}
  \label{fig:tl-summary}
\end{figure}

\draftnote{state the smallest displacement at which the difference image in
\cref{fig:tl-summary}b still shows a clear, unambiguous peak above the
numerical background, and compare it with the amplitude resolution floor
found for stationary scatterers in \cref{sec:meth-resolution}. The noisy
version of this sweep, and the noise-mitigation strategies tested against
it, are presented in \cref{ch:hyp3} rather than here.}

\section{Vertical Movement}
\label{sec:hyp1-vertical}

This section mirrors \cref{sec:hyp1-lateral}, replacing lateral displacement
of the scatterer with \emph{vertical} (depth) displacement. The two
directions are not expected to behave identically: \cref{sec:th-duality}
showed that a migrated image separates lateral spatial frequencies like a
prism but does not separate vertical ones, so the amplitude-based
detectability established here is an important point of comparison for the
phase-based vertical results of \cref{sec:hyp1-phaseplane}.

A single PEC cylinder ($r=\SI{28}{mm}$) sits at lateral position
$x=\SI{2.0}{m}$ and is displaced \emph{downward} from its baseline depth
across seven scenarios, from $1\lambda$ down to $\tfrac{1}{32}\lambda$
(\cref{fig:vtl-setup}), using the same domain and grid as
\cref{sec:meth-resolution,sec:hyp1-lateral}.

\begin{figure}[htbp]
  \centering
  \begin{subfigure}[t]{0.48\textwidth}
    \includegraphics[width=\linewidth]{VTL_001_Vertical_TimeLapse_Study__Model_Geometry__domain_4010_m_Δx.png}
    \caption{Model geometry}
  \end{subfigure}
  \hfill
  \begin{subfigure}[t]{0.48\textwidth}
    \includegraphics[width=\linewidth]{VTL_002_Moving_scatterer__all_7_scenarios__r__28_mm_x__20_mBaseline.png}
    \caption{All 7 vertical displacement scenarios}
  \end{subfigure}
  \caption{Forward-model setup for the vertical time-lapse study: (a) the
    gprMax domain and grid; (b) the vertical displacement of the scatterer
    across all seven scenarios, $r=\SI{28}{mm}$, fixed at $x=\SI{2.0}{m}$.}
  \label{fig:vtl-setup}
\end{figure}

\Cref{fig:vtl-bscans} shows the background-baseline and vertically
time-lapsed raw B-scans and the background-subtracted result, with the
expected arrival times for each scenario marked, and \cref{fig:vtl-taper}
the effect of the standard tapering and $t_0$-shift conditioning on the
$1\lambda$ dataset.

\begin{figure}[htbp]
  \centering
  \begin{subfigure}[t]{0.48\textwidth}
    \includegraphics[width=\linewidth]{VTL_003_GPR_B-Scans__Background_Baseline_and_Vertical_TimeLapsed_Mod.png}
    \caption{Background \& time-lapsed}
  \end{subfigure}
  \hfill
  \begin{subfigure}[t]{0.48\textwidth}
    \includegraphics[width=\linewidth]{VTL_004_GPR_B-Scans__Background_Subtracted__green_dashed__expected_a.png}
    \caption{Background-subtracted}
  \end{subfigure}
  \caption{Raw B-scans for the vertical time-lapse study: (a) background and
    time-lapsed models; (b) background-subtracted, with expected arrival
    times marked (green, dashed).}
  \label{fig:vtl-bscans}
\end{figure}

\begin{figure}[htbp]
  \centering
  \begin{subfigure}[t]{0.48\textwidth}
    \includegraphics[width=\linewidth]{VTL_005_Effect_of_Tapering_and_t0_Shift__1λ_dataset_single_trace.png}
    \caption{Single trace}
  \end{subfigure}
  \hfill
  \begin{subfigure}[t]{0.48\textwidth}
    \includegraphics[width=\linewidth]{VTL_006_B-scan_effect_of_tapering_and_t0_shift__1λ_dataset.png}
    \caption{Full B-scan}
  \end{subfigure}
  \caption{Effect of tapering and the $t_0$ shift on the $1\lambda$ vertical
    displacement dataset: (a) a single trace; (b) the full B-scan.}
  \label{fig:vtl-taper}
\end{figure}

\Cref{fig:vtl-kirchhoff,fig:vtl-gazdag,fig:vtl-backprop} repeat the same
three-algorithm migration-and-differencing analysis for the vertical
displacement scenarios.

\begin{figure}[htbp]
  \centering
  \begin{subfigure}[t]{0.48\textwidth}
    \includegraphics[width=\linewidth]{VTL_007_Kirchhoff_Migration__All_7_Datasets____f_c15_GHz____aperture.png}
    \caption{All 7 scenarios}
  \end{subfigure}
  \hfill
  \begin{subfigure}[t]{0.48\textwidth}
    \includegraphics[width=\linewidth]{VTL_008_Kirchhoff_Migration_zoomed____f_c15_GHz____aperture40.png}
    \caption{Zoomed}
  \end{subfigure}

  \vspace{0.6em}
  \begin{subfigure}[t]{0.48\textwidth}
    \includegraphics[width=\linewidth]{VTL_009_Kirchhoff_Migration__TimeLapse_Differences_migrated__migrate.png}
    \caption{Time-lapse difference}
  \end{subfigure}
  \hfill
  \begin{subfigure}[t]{0.48\textwidth}
    \includegraphics[width=\linewidth]{VTL_010_Kirchhoff_Migration__TimeLapse_Differences_zoomed.png}
    \caption{Difference, zoomed}
  \end{subfigure}
  \caption{Kirchhoff migration of the vertical time-lapse study
    ($f_c=\SI{15}{\GHz}$, aperture $=40$): (a,b) migrated image for all
    scenarios and zoomed; (c,d) time-lapse difference and zoomed.}
  \label{fig:vtl-kirchhoff}
\end{figure}

\begin{figure}[htbp]
  \centering
  \begin{subfigure}[t]{0.48\textwidth}
    \includegraphics[width=\linewidth]{VTL_011_Gazdag_Phase-Shift_Migration__All_7_Datasets____f_c15_GHz.png}
    \caption{All scenarios}
  \end{subfigure}
  \hfill
  \begin{subfigure}[t]{0.48\textwidth}
    \includegraphics[width=\linewidth]{VTL_012_Gazdag_Phase-Shift_Migration_zoomed____f_cf_c_GHz_GHz.png}
    \caption{Zoomed}
  \end{subfigure}

  \vspace{0.6em}
  \begin{subfigure}[t]{0.48\textwidth}
    \includegraphics[width=\linewidth]{VTL_013_Gazdag_Migration__TimeLapse_Differences_migrated__migrated_b.png}
    \caption{Time-lapse difference}
  \end{subfigure}
  \hfill
  \begin{subfigure}[t]{0.48\textwidth}
    \includegraphics[width=\linewidth]{VTL_014_Gazdag_Migration__TimeLapse_Differences_zoomed.png}
    \caption{Difference, zoomed}
  \end{subfigure}
  \caption{Gazdag phase-shift migration of the vertical time-lapse study
    ($f_c=\SI{15}{\GHz}$): (a,b) migrated image and zoomed; (c,d) time-lapse
    difference and zoomed.}
  \label{fig:vtl-gazdag}
\end{figure}

\begin{figure}[htbp]
  \centering
  \begin{subfigure}[t]{0.48\textwidth}
    \includegraphics[width=\linewidth]{VTL_015_Back-Propagation_E__All_7_Datasets____focus_at_1906_ns.png}
    \caption{$\lVert\mathbf{E}\rVert$, all scenarios}
  \end{subfigure}
  \hfill
  \begin{subfigure}[t]{0.48\textwidth}
    \includegraphics[width=\linewidth]{VTL_016_Back-Propagation_E_zoomed____focus_at_1906_ns.png}
    \caption{$\lVert\mathbf{E}\rVert$, zoomed}
  \end{subfigure}

  \vspace{0.6em}
  \begin{subfigure}[t]{0.48\textwidth}
    \includegraphics[width=\linewidth]{VTL_017_Back-Propagation_Ez__All_7_Datasets____focus_at_1906_ns.png}
    \caption{$E_z$, all scenarios}
  \end{subfigure}
  \hfill
  \begin{subfigure}[t]{0.48\textwidth}
    \includegraphics[width=\linewidth]{VTL_018_Back-Propagation_Ez_zoomed____focus_at_1906_ns.png}
    \caption{$E_z$, zoomed}
  \end{subfigure}

  \vspace{0.6em}
  \begin{subfigure}[t]{0.48\textwidth}
    \includegraphics[width=\linewidth]{VTL_019_Back-Propagation__TimeLapse_Differences_Ez_Ez__Ez_baseline.png}
    \caption{$E_z$ time-lapse difference}
  \end{subfigure}
  \hfill
  \begin{subfigure}[t]{0.48\textwidth}
    \includegraphics[width=\linewidth]{VTL_020_Back-Propagation__TimeLapse_Differences_Ez_zoomed.png}
    \caption{Difference, zoomed}
  \end{subfigure}
  \caption{Time-reversal back-propagation migration of the vertical
    time-lapse study, focused at $t=\SI{1906}{\ns}$: field-magnitude (a,b)
    and $E_z$ (c,d) images, and the $E_z$ time-lapse difference (e,f).}
  \label{fig:vtl-backprop}
\end{figure}

\subsection*{Vertical Detectability Summary}
\label{sec:vtl-detectability}

\Cref{fig:vtl-summary} overlays the signed time-lapse-difference amplitude
from all three algorithms and plots the normalised \emph{vertical}
point-spread function of that difference at $x=\SI{2.0}{m}$, as a function
of vertical displacement.

\begin{figure}[htbp]
  \centering
  \begin{subfigure}[t]{0.48\textwidth}
    \includegraphics[width=\linewidth]{VTL_021_Vertical_TimeLapse_Migration_Comparison__Signed_Amplitude.png}
    \caption{Signed-amplitude comparison}
  \end{subfigure}
  \hfill
  \begin{subfigure}[t]{0.48\textwidth}
    \includegraphics[width=\linewidth]{VTL_022_Normalised_Vertical_PSF__TimeLapse_Difference_at_x__20_m.png}
    \caption{Normalised vertical PSF of the difference}
  \end{subfigure}
  \caption{(a) Signed time-lapse-difference amplitude for all three
    migration algorithms; (b) the normalised vertical PSF of the difference
    image at $x=\SI{2.0}{m}$, swept across vertical displacements from
    $1\lambda$ to $\tfrac{1}{32}\lambda$.}
  \label{fig:vtl-summary}
\end{figure}

\draftnote{compare the smallest reliably-detected vertical displacement in
\cref{fig:vtl-summary}b against the lateral result of
\cref{fig:tl-summary}b (\cref{sec:tl-detectability}) and discuss whether the
lateral/vertical asymmetry predicted by the instantaneous-phase theory of
\cref{sec:th-duality} is also visible at the amplitude level, or only
emerges once the phase-plane estimator of \cref{sec:hyp1-phaseplane} is
applied.}

\section{Diagonal Movement}
\label{sec:hyp1-diagonal}

The third and final translation direction combines the previous two: the
scatterer moves simultaneously laterally and vertically, along a fixed
$2{:}1$ diagonal ($\Dx = 2\,\Dz$ in every scenario), so that all scenarios
lie on the same line through the baseline position. A single PEC cylinder
at $x=\SI{2.0}{m}$, baseline depth \SI{0.676}{m}, is displaced diagonally
across five scenarios (\cref{tab:dtl-scenarios}), using the same domain and
grid as \cref{sec:meth-resolution,sec:hyp1-lateral,sec:hyp1-vertical}.

\begin{table}[htbp]
  \centering
  \caption{Diagonal displacement scenarios: lateral ($\Dx$, rightward) and
    vertical ($\Dz$, downward) components, both expressed as fractions of
    the dominant wavelength $\lambda$.}
  \label{tab:dtl-scenarios}
  \begin{tabular}{lcc}
    \toprule
    Scenario & $\Dx$ (right) & $\Dz$ (down) \\
    \midrule
    Baseline   & 0                       & 0 \\
    1          & $1\lambda$              & $\tfrac12\lambda$ \\
    2          & $\tfrac12\lambda$       & $\tfrac14\lambda$ \\
    3          & $\tfrac14\lambda$       & $\tfrac18\lambda$ \\
    4          & $\tfrac18\lambda$       & $\tfrac{1}{16}\lambda$ \\
    5          & $\tfrac{1}{16}\lambda$  & $\tfrac{1}{32}\lambda$ \\
    \bottomrule
  \end{tabular}
\end{table}

\begin{figure}[htbp]
  \centering
  \begin{subfigure}[t]{0.48\textwidth}
    \includegraphics[width=\linewidth]{DTL_001_Diagonal_TimeLapse_Study__Model_Geometry__domain_4010_m_Δx.png}
    \caption{Model geometry}
  \end{subfigure}
  \hfill
  \begin{subfigure}[t]{0.48\textwidth}
    \includegraphics[width=\linewidth]{DTL_002_Diagonal_scatterer_path__all_6_scenarios__r__28_mm.png}
    \caption{All scenarios, diagonal path}
  \end{subfigure}
  \caption{Forward-model setup for the diagonal time-lapse study: (a) the
    gprMax domain and grid; (b) the diagonal displacement path of the
    scatterer across all scenarios, $r=\SI{28}{mm}$.}
  \label{fig:dtl-setup}
\end{figure}

\Cref{fig:dtl-bscans} shows the background-baseline and diagonally
time-lapsed raw B-scans, the background-subtracted result, and the effect
of the synthetic Laplace noise used later in \cref{ch:hyp3}; \cref{fig:dtl-taper}
shows the standard tapering and $t_0$-shift conditioning on the scenario-1
dataset.

\begin{figure}[htbp]
  \centering
  \begin{subfigure}[t]{0.31\textwidth}
    \includegraphics[width=\linewidth]{DTL_003_GPR_B-Scans__Background_Baseline_and_Diagonal_TimeLapsed_Mod.png}
    \caption{Background \& time-lapsed}
  \end{subfigure}
  \hfill
  \begin{subfigure}[t]{0.31\textwidth}
    \includegraphics[width=\linewidth]{DTL_004_GPR_B-Scans__Background_Subtracted__green_dashed__expected_a.png}
    \caption{Background-subtracted}
  \end{subfigure}
  \hfill
  \begin{subfigure}[t]{0.31\textwidth}
    \includegraphics[width=\linewidth]{DTL_005_GPR_B-Scans__With_Synthetic_Laplace_Noise_10_of_signal_std.png}
    \caption{With Laplace noise}
  \end{subfigure}
  \caption{Raw B-scans for the diagonal time-lapse study: (a) background and
    time-lapsed models; (b) background-subtracted, with expected arrival
    times marked; (c) with synthetic Laplace-distributed noise at $10\,\%$
    of the signal standard deviation (used in \cref{ch:hyp3}).}
  \label{fig:dtl-bscans}
\end{figure}

\begin{figure}[htbp]
  \centering
  \begin{subfigure}[t]{0.48\textwidth}
    \includegraphics[width=\linewidth]{DTL_007_Effect_of_Tapering_and_t0_Shift__Scenario_1_dataset_single_t.png}
    \caption{Single trace}
  \end{subfigure}
  \hfill
  \begin{subfigure}[t]{0.48\textwidth}
    \includegraphics[width=\linewidth]{DTL_008_B-scan_effect_of_tapering_and_t0_shift__Scenario_1_dataset.png}
    \caption{Full B-scan}
  \end{subfigure}
  \caption{Effect of tapering and the $t_0$ shift on the scenario-1
    ($1\lambda,\tfrac12\lambda$) diagonal displacement dataset: (a) a single
    trace; (b) the full B-scan.}
  \label{fig:dtl-taper}
\end{figure}

\Cref{fig:dtl-kirchhoff,fig:dtl-gazdag,fig:dtl-backprop} repeat the same
three-algorithm migration-and-differencing analysis for the diagonal
displacement scenarios.

\begin{figure}[htbp]
  \centering
  \begin{subfigure}[t]{0.48\textwidth}
    \includegraphics[width=\linewidth]{DTL_009_Kirchhoff_Migration__All_6_Datasets____f_c15_GHz____aperture.png}
    \caption{All 6 scenarios}
  \end{subfigure}
  \hfill
  \begin{subfigure}[t]{0.48\textwidth}
    \includegraphics[width=\linewidth]{DTL_010_Kirchhoff_Migration_zoomed____f_c15_GHz____aperture40.png}
    \caption{Zoomed}
  \end{subfigure}

  \vspace{0.6em}
  \begin{subfigure}[t]{0.48\textwidth}
    \includegraphics[width=\linewidth]{DTL_011_Kirchhoff_Migration__TimeLapse_Differences_migrated__migrate.png}
    \caption{Time-lapse difference}
  \end{subfigure}
  \hfill
  \begin{subfigure}[t]{0.48\textwidth}
    \includegraphics[width=\linewidth]{DTL_012_Kirchhoff_Migration__TimeLapse_Differences_zoomed.png}
    \caption{Difference, zoomed}
  \end{subfigure}
  \caption{Kirchhoff migration of the diagonal time-lapse study
    ($f_c=\SI{15}{\GHz}$, aperture $=40$): (a,b) migrated image for all
    scenarios and zoomed; (c,d) time-lapse difference and zoomed.}
  \label{fig:dtl-kirchhoff}
\end{figure}

\begin{figure}[htbp]
  \centering
  \begin{subfigure}[t]{0.48\textwidth}
    \includegraphics[width=\linewidth]{DTL_013_Gazdag_Phase-Shift_Migration__All_6_Datasets____f_c15_GHz.png}
    \caption{All scenarios}
  \end{subfigure}
  \hfill
  \begin{subfigure}[t]{0.48\textwidth}
    \includegraphics[width=\linewidth]{DTL_014_Gazdag_Phase-Shift_Migration_zoomed____f_c15_GHz.png}
    \caption{Zoomed}
  \end{subfigure}

  \vspace{0.6em}
  \begin{subfigure}[t]{0.48\textwidth}
    \includegraphics[width=\linewidth]{DTL_015_Gazdag_Migration__TimeLapse_Differences_migrated__migrated_b.png}
    \caption{Time-lapse difference}
  \end{subfigure}
  \hfill
  \begin{subfigure}[t]{0.48\textwidth}
    \includegraphics[width=\linewidth]{DTL_016_Gazdag_Migration__TimeLapse_Differences_zoomed.png}
    \caption{Difference, zoomed}
  \end{subfigure}
  \caption{Gazdag phase-shift migration of the diagonal time-lapse study
    ($f_c=\SI{15}{\GHz}$): (a,b) migrated image and zoomed; (c,d) time-lapse
    difference and zoomed.}
  \label{fig:dtl-gazdag}
\end{figure}

\begin{figure}[htbp]
  \centering
  \begin{subfigure}[t]{0.48\textwidth}
    \includegraphics[width=\linewidth]{DTL_017_Back-Propagation_E__All_6_Datasets____focus_at_1906_ns.png}
    \caption{$\lVert\mathbf{E}\rVert$, all scenarios}
  \end{subfigure}
  \hfill
  \begin{subfigure}[t]{0.48\textwidth}
    \includegraphics[width=\linewidth]{DTL_018_Back-Propagation_E_zoomed____focus_at_1906_ns.png}
    \caption{$\lVert\mathbf{E}\rVert$, zoomed}
  \end{subfigure}

  \vspace{0.6em}
  \begin{subfigure}[t]{0.48\textwidth}
    \includegraphics[width=\linewidth]{DTL_019_Back-Propagation_Ez__All_6_Datasets____focus_at_1906_ns.png}
    \caption{$E_z$, all scenarios}
  \end{subfigure}
  \hfill
  \begin{subfigure}[t]{0.48\textwidth}
    \includegraphics[width=\linewidth]{DTL_020_Back-Propagation_Ez_zoomed____focus_at_1906_ns.png}
    \caption{$E_z$, zoomed}
  \end{subfigure}

  \vspace{0.6em}
  \begin{subfigure}[t]{0.48\textwidth}
    \includegraphics[width=\linewidth]{DTL_021_Back-Propagation__TimeLapse_Differences_Ez_Ez__Ez_baseline.png}
    \caption{$E_z$ time-lapse difference}
  \end{subfigure}
  \hfill
  \begin{subfigure}[t]{0.48\textwidth}
    \includegraphics[width=\linewidth]{DTL_022_Back-Propagation__TimeLapse_Differences_Ez_zoomed.png}
    \caption{Difference, zoomed}
  \end{subfigure}
  \caption{Time-reversal back-propagation migration of the diagonal
    time-lapse study, focused at $t=\SI{1906}{\ns}$: field-magnitude (a,b)
    and $E_z$ (c,d) images, and the $E_z$ time-lapse difference (e,f).}
  \label{fig:dtl-backprop}
\end{figure}

\subsection*{Diagonal Detectability Summary}
\label{sec:dtl-detectability}

Because every scenario moves the scatterer along the same $2{:}1$ diagonal
line through the baseline, the detectability analysis samples each
time-lapse-difference image \emph{along that line} rather than along a
single Cartesian axis: \cref{fig:dtl-summary}a overlays the signed,
diagonally-sampled amplitude from all three algorithms, and
\cref{fig:dtl-summary}b plots the normalised diagonal point-spread function
(signed amplitude and Hilbert envelope) as a function of signed distance
along the motion direction.

\begin{figure}[htbp]
  \centering
  \begin{subfigure}[t]{0.48\textwidth}
    \includegraphics[width=\linewidth]{DTL_023_Diagonal_TimeLapse_Migration_Comparison__Signed_Amplitude.png}
    \caption{Signed-amplitude comparison}
  \end{subfigure}
  \hfill
  \begin{subfigure}[t]{0.48\textwidth}
    \includegraphics[width=\linewidth]{DTL_024_Normalised_Diagonal_PSF__TimeLapse_Difference_Along_Motion_D.png}
    \caption{Normalised diagonal PSF along motion direction}
  \end{subfigure}
  \caption{(a) Signed time-lapse-difference amplitude for all three
    migration algorithms, sampled along the diagonal motion direction;
    (b) the normalised diagonal PSF of the difference image along that same
    direction, swept across the five diagonal scenarios of
    \cref{tab:dtl-scenarios}.}
  \label{fig:dtl-summary}
\end{figure}

\draftnote{state the smallest diagonal displacement at which
\cref{fig:dtl-summary}b still shows a clear, unambiguous peak above
background, and compare it with the lateral and vertical floors of
\cref{sec:tl-detectability,sec:vtl-detectability} --- since the diagonal
step combines a lateral and a vertical component of different magnitude
($\Dx=2\Dz$), state whether the effective detectability floor tracks the
(tighter) lateral floor, the (looser) vertical floor, or some combination of
the two.}

\section{Global Phase-Plane Validation}
\label{sec:hyp1-phaseplane}

\Cref{sec:hyp1-lateral,sec:hyp1-vertical,sec:hyp1-diagonal} showed that
simple amplitude differencing of migrated images detects sub-wavelength
displacement only down to a finite floor, in all three directions. This
section applies the phase-plane shift estimator derived in
\cref{sec:th-fourier-shift,sec:th-wls,sec:th-material-change} and
implemented as described in \cref{sec:meth-phaseplane} to the lateral and
vertical displacement datasets, and shows that it remains accurate at
displacement scales where the amplitude image is already featureless.

Because several of the analyses below are repeated identically across seven
(or six) wavelength scales, only one or two representative scales are shown
in this section; the complete sweep for every analysis is given in
\cref{app:extended-sweeps}, so that every one of the relevant figures
produced by \texttt{TimeLapse\_Processing.ipynb} for these scenarios appears
at least once in the thesis.

\subsection*{Migrated Baseline Overview Across Scenarios}
\label{sec:tlp-overview}

\Cref{fig:tlp-overview} shows the starting point for every analysis in this
section: the raw background-subtracted B-scans and the three migrated
images (Kirchhoff, Gazdag, back-propagation) for every scenario used below.

\begin{figure}[htbp]
  \centering
  \begin{subfigure}[t]{0.48\textwidth}
    \includegraphics[width=\linewidth]{TLP_001_Raw_Background-Subtracted_B-Scans__All_Scenarios.png}
    \caption{Raw, background-subtracted}
  \end{subfigure}
  \hfill
  \begin{subfigure}[t]{0.48\textwidth}
    \includegraphics[width=\linewidth]{TLP_002_Kirchhoff_Migration__All_Scenarios.png}
    \caption{Kirchhoff-migrated}
  \end{subfigure}

  \vspace{0.6em}
  \begin{subfigure}[t]{0.48\textwidth}
    \includegraphics[width=\linewidth]{TLP_003_Gazdag_Migration__All_Scenarios.png}
    \caption{Gazdag-migrated}
  \end{subfigure}
  \hfill
  \begin{subfigure}[t]{0.48\textwidth}
    \includegraphics[width=\linewidth]{TLP_004_Back-prop_Migration__All_Scenarios.png}
    \caption{Back-propagation-migrated}
  \end{subfigure}
  \caption{Overview of all scenarios re-used from \cref{sec:hyp1-lateral} in
    this section: (a) raw background-subtracted B-scans; (b--d) the same
    data migrated with Kirchhoff, Gazdag, and back-propagation migration.}
  \label{fig:tlp-overview}
\end{figure}

\subsection*{Phase-Plane Fit: Lateral Shift Validation}
\label{sec:tlp-horizontal-validation}

For every lateral-displacement scenario of \cref{sec:hyp1-lateral}, the
estimator of \cref{sec:meth-phaseplane} is applied to the baseline/monitor
pair, cropped to $\pm2.5\lambda$ around the scatterer apex (located from the
Hilbert envelope of the baseline image). \Cref{fig:tlp-horiz-validation}
shows the recovered $(\Dz,\Dx)$ compared with the known ground truth for all
three migration methods.

\begin{figure}[htbp]
  \centering
  \begin{subfigure}[t]{0.32\textwidth}
    \includegraphics[width=\linewidth]{TLP_005_Phase-plane_shift_estimation__Kirchhoff____Baseline_vs_each.png}
    \caption{Kirchhoff}
  \end{subfigure}
  \hfill
  \begin{subfigure}[t]{0.32\textwidth}
    \includegraphics[width=\linewidth]{TLP_006_Phase-plane_shift_estimation__Gazdag____Baseline_vs_each_sce.png}
    \caption{Gazdag}
  \end{subfigure}
  \hfill
  \begin{subfigure}[t]{0.32\textwidth}
    \includegraphics[width=\linewidth]{TLP_007_Phase-plane_shift_estimation__Back-prop____Baseline_vs_each.png}
    \caption{Back-propagation}
  \end{subfigure}
  \caption{2D phase-plane shift estimation applied to the baseline-versus-each-scenario
    pairs of \cref{sec:hyp1-lateral}, for all three migration methods.}
  \label{fig:tlp-horiz-validation}
\end{figure}

\draftnote{read off \cref{fig:tlp-horiz-validation} the smallest lateral
displacement at which the estimated $\Dx$ still tracks the true value to
within (state your accepted tolerance, e.g.\ $\pm 5\,\%$ or $\pm\SI{1}{mm}$),
for each migration method, and contrast this explicitly with the amplitude
floor found in \cref{sec:tl-detectability}.}

\subsection*{Phase-Plane Fit: Vertical Shift Validation}
\label{sec:tlp-vertical}

\Cref{fig:tlp-vert-validation} repeats the previous analysis for the
vertical-displacement dataset of \cref{sec:hyp1-vertical}, with the ROI
crop and cross-spectrum fit unchanged except that the displacement being
recovered is now $\Dz$ rather than $\Dx$.

\begin{figure}[htbp]
  \centering
  \begin{subfigure}[t]{0.32\textwidth}
    \includegraphics[width=\linewidth]{TLP_010_VerticalTimeLapse__Phase-plane_shift_estimation__Kirchhoff.png}
    \caption{Kirchhoff}
  \end{subfigure}
  \hfill
  \begin{subfigure}[t]{0.32\textwidth}
    \includegraphics[width=\linewidth]{TLP_011_VerticalTimeLapse__Phase-plane_shift_estimation__Gazdag____B.png}
    \caption{Gazdag}
  \end{subfigure}
  \hfill
  \begin{subfigure}[t]{0.32\textwidth}
    \includegraphics[width=\linewidth]{TLP_012_VerticalTimeLapse__Phase-plane_shift_estimation__Back-prop.png}
    \caption{Back-propagation}
  \end{subfigure}
  \caption{2D phase-plane shift estimation applied to the baseline-versus-each-scenario
    pairs of \cref{sec:hyp1-vertical}, for all three migration methods.}
  \label{fig:tlp-vert-validation}
\end{figure}

\draftnote{state whether the vertical phase-plane fit remains accurate to
the same smallest displacement found for the lateral case in
\cref{sec:tlp-horizontal-validation}, or whether the lateral/vertical
asymmetry predicted in \cref{sec:th-duality} (a constant vertical
\emph{plateau} versus a lateral \emph{slope}) shows up as a difference in
achievable precision between \cref{fig:tlp-horiz-validation,fig:tlp-vert-validation}.}

\subsection*{Phase-Plane Fit: Diagonal Shift Validation}
\label{sec:tlp-diagonal}

\draftnote{\textbf{Gap:} unlike the lateral and vertical cases above, the 2D
WLS phase-plane fit has not yet been run on the diagonal dataset of
\cref{sec:hyp1-diagonal} --- \texttt{Diagonal\_TimeLapse\_Playground.ipynb}
currently only produces the amplitude-based detectability result of
\cref{sec:dtl-detectability}. Completing Hypothesis 1's diagonal claim
requires applying \texttt{estimate\_shift\_2d} (\cref{sec:meth-phaseplane})
to the same baseline/monitor pairs and recovering the joint $(\Dz,\Dx)$
estimate, which should satisfy the known $\Dx=2\Dz$ ratio of
\cref{tab:dtl-scenarios} if the fit is working correctly --- this is the
single most important piece of evidence still missing from this chapter.}

\section{Hypothesis 1.5 (Optional): Local Phase-Gradient Methods}
\label{sec:hyp1-h15}

The phase-plane fit of \cref{sec:hyp1-phaseplane} is a \emph{global}
operation: it acts on the whole 2D spectrum at once. This section explores
a \emph{local} alternative --- the spatial instantaneous phase, extracted
trace-by-trace with a Riesz/Hilbert transform, and its time-frequency
relatives --- to test Hypothesis 1.5: whether the local phase gradients
$\partial\varphi/\partial x,\partial\varphi/\partial y$ offer an equivalent,
simpler route to the same displacement estimate. \draftnote{this section is
exploratory: the underlying Riesz-transform/monogenic-signal machinery is
implemented in \texttt{CWT\_playground.ipynb}, but it has not yet been
validated quantitatively against the global WLS fit of
\cref{sec:hyp1-phaseplane} on the same datasets --- treat the results below
as a qualitative complement to, not a replacement for, Hypothesis 1's
primary evidence.}

\subsection{Instantaneous Phase Imaging --- Lateral Displacement}
\label{sec:tlp-instphase-horizontal}

The 2D Riesz transform of \cref{eq:kx-inst}--\cref{eq:dphi-lateral} is
applied directly to the Gazdag-migrated baseline and monitor images,
producing an instantaneous-phase difference map $\Dphi(z,x)$ and, at the
scatterer depth, a 1D phase cross-section whose slope at the zero-crossing
encodes $\Dx$. This analysis is repeated at all seven lateral scales from
$2\lambda$ to $\tfrac{1}{32}\lambda$; \cref{fig:tlp-instphase-horiz-example}
shows the representative $\tfrac14\lambda$ case ($\Dx=\SI{280}{mm}$), and
\cref{fig:tlp-instphase-horiz-summary} summarises the fitted zero-crossing
slope across all seven scales. The full seven-scale sweep is given in
\cref{app:sec-instphase-horizontal}.

\begin{figure}[htbp]
  \centering
  \begin{subfigure}[t]{0.32\textwidth}
    \includegraphics[width=\linewidth]{TLP_025_Instantaneous_Phase__Gazdag____Baseline_vs_¼λ___Δx__280_mm.png}
    \caption{Instantaneous phase}
  \end{subfigure}
  \hfill
  \begin{subfigure}[t]{0.32\textwidth}
    \includegraphics[width=\linewidth]{TLP_026_Phase_cross-section__z__676_mm____Gazdag____¼λ__Δx__280_mm.png}
    \caption{Cross-section, $z=\SI{676}{mm}$}
  \end{subfigure}
  \hfill
  \begin{subfigure}[t]{0.32\textwidth}
    \includegraphics[width=\linewidth]{TLP_027_Phase_cross-section__zoomed__½λ_around_scatterers.png}
    \caption{Cross-section, zoomed}
  \end{subfigure}
  \caption{Instantaneous-phase analysis for the representative
    $\Dx=\tfrac14\lambda=\SI{280}{mm}$ lateral displacement, Gazdag-migrated:
    (a) the 2D wrapped phase-difference map; (b) the 1D phase cross-section
    at the scatterer depth $z=\SI{676}{mm}$; (c) the same cross-section
    zoomed to $\pm\tfrac12\lambda$ around the scatterer.}
  \label{fig:tlp-instphase-horiz-example}
\end{figure}

\begin{figure}[htbp]
  \centering
  \includegraphics[width=0.7\linewidth]{TLP_037_Gazdag__Δφ_zero-crossing_slope_vs_scatterer_separation.png}
  \caption{Fitted zero-crossing slope of the lateral instantaneous
    phase-difference cross-section, as a function of true scatterer
    displacement, across all seven scales from $2\lambda$ to
    $\tfrac{1}{32}\lambda$ (Gazdag-migrated).}
  \label{fig:tlp-instphase-horiz-summary}
\end{figure}

\subsection{Instantaneous Phase Imaging --- Vertical Displacement}
\label{sec:tlp-instphase-vertical}

The same instantaneous-phase analysis is repeated for the vertical
displacement dataset, with the cross-section now taken along $z$ at the
fixed lateral position $x=\SI{2.0}{m}$ of the scatterer, across six scales
from $1\lambda$ to $\tfrac{1}{32}\lambda$. \Cref{fig:tlp-instphase-vert-example}
shows the representative $\tfrac14\lambda$ case ($\Dz=\SI{280}{mm}$), and
\cref{fig:tlp-instphase-vert-summary} summarises the mean phase difference
across all six scales. The full sweep is given in
\cref{app:sec-instphase-vertical}.

\begin{figure}[htbp]
  \centering
  \begin{subfigure}[t]{0.32\textwidth}
    \includegraphics[width=\linewidth]{TLP_044_VerticalTimeLapse__Instantaneous_Phase__Gazdag____Baseline_v.png}
    \caption{Instantaneous phase}
  \end{subfigure}
  \hfill
  \begin{subfigure}[t]{0.32\textwidth}
    \includegraphics[width=\linewidth]{TLP_045_Phase_cross-section__x__2000_mm____Gazdag____¼λ__Δz__280_mm.png}
    \caption{Cross-section, $x=\SI{2000}{mm}$}
  \end{subfigure}
  \hfill
  \begin{subfigure}[t]{0.32\textwidth}
    \includegraphics[width=\linewidth]{TLP_046_Phase_cross-section__zoomed__½λ_around_scatterer_depths.png}
    \caption{Cross-section, zoomed}
  \end{subfigure}
  \caption{Instantaneous-phase analysis for the representative
    $\Dz=\tfrac14\lambda=\SI{280}{mm}$ vertical displacement, Gazdag-migrated:
    (a) the 2D wrapped phase-difference map; (b) the 1D phase cross-section
    at $x=\SI{2000}{mm}$; (c) the same cross-section zoomed to
    $\pm\tfrac12\lambda$ around the scatterer depths.}
  \label{fig:tlp-instphase-vert-example}
\end{figure}

\begin{figure}[htbp]
  \centering
  \includegraphics[width=0.7\linewidth]{TLP_056_VerticalTimeLapse__Gazdag__Mean_Δφ_vs_scatterer_vertical_shi.png}
  \caption{Mean instantaneous phase difference $\Dphi$ as a function of true
    vertical scatterer shift, across all six scales from $1\lambda$ to
    $\tfrac{1}{32}\lambda$ (Gazdag-migrated). Unlike the lateral case
    (\cref{fig:tlp-instphase-horiz-summary}), this is a plateau level rather
    than a cross-section slope, consistent with the asymmetry derived in
    \cref{sec:th-duality}.}
  \label{fig:tlp-instphase-vert-summary}
\end{figure}

\subsection{Spectral-Line Phase Analysis}
\label{sec:tlp-spectral-line}

The localised-Fourier-shift spectral-line equation
\cref{eq:local-phase-line} is verified directly on individual traces using
the CLSSA decomposition of \texttt{phase\_decomposition.py}, applied both to
the Gazdag-migrated image and to the raw, unmigrated B-scan, at the column
position containing each scenario's scatterer. \Cref{fig:tlp-spectral-line-example}
compares the representative $\tfrac14\lambda$ case in both domains; the full
seven-scale sweep for both is given in
\cref{app:sec-spectral-line}.

\begin{figure}[htbp]
  \centering
  \begin{subfigure}[t]{0.48\textwidth}
    \includegraphics[width=\linewidth]{TLP_060_Spectral_Line_CLSSA____Gazdag____¼λ___Δx__280_mm__02500λ.png}
    \caption{Gazdag-migrated}
  \end{subfigure}
  \hfill
  \begin{subfigure}[t]{0.48\textwidth}
    \includegraphics[width=\linewidth]{TLP_067_Spectral_Line_CLSSA____Raw_Unmigrated____¼λ___Δx__280_mm__02.png}
    \caption{Raw, unmigrated}
  \end{subfigure}
  \caption{CLSSA spectral-line decomposition (amplitude spectrum, phase
    spectrum, and phase gather, \cref{sec:th-local}) at the representative
    $\Dx=\tfrac14\lambda=\SI{280}{mm}$ scale: (a) Gazdag-migrated trace; (b)
    raw, unmigrated trace at the same scenario.}
  \label{fig:tlp-spectral-line-example}
\end{figure}

\draftnote{state whether the migrated and raw spectral lines in
\cref{fig:tlp-spectral-line-example} give consistent $\Dt$/slope estimates
once converted with \cref{eq:dt-dz}, which would confirm that migration is
not required for the temporal spectral-line method to work, only for the
2D wavenumber plane fit of \cref{sec:th-wls}.}

\subsection{Cross-Phase Spectrograms}
\label{sec:tlp-spectrogram}

\Cref{fig:tlp-spectrogram-example} shows the cross-phase spectrogram
$\Dphi(\tau,f)$ of \cref{sec:th-local} for two contrasting scales: a large,
easily visible $\tfrac14\lambda$ displacement and the smallest,
$\tfrac{1}{32}\lambda$, where the coherent window around the scatterer's
two-way time becomes much harder to distinguish from the incoherent
background. The full seven-scale sweep is given in
\cref{app:sec-spectrogram}.

\begin{figure}[htbp]
  \centering
  \begin{subfigure}[t]{0.48\textwidth}
    \includegraphics[width=\linewidth]{TLP_074_Cross-Phase_Spectrogram_CLSSA__ΔΦτf____Gazdag____¼λ.png}
    \caption{$\Dx=\tfrac14\lambda$}
  \end{subfigure}
  \hfill
  \begin{subfigure}[t]{0.48\textwidth}
    \includegraphics[width=\linewidth]{TLP_077_Cross-Phase_Spectrogram_CLSSA__ΔΦτf____Gazdag____¹₃₂λ.png}
    \caption{$\Dx=\tfrac{1}{32}\lambda$}
  \end{subfigure}
  \caption{Cross-phase spectrogram $\Dphi(\tau,f)$ (two-way time against
    frequency) at two contrasting lateral displacement scales,
    Gazdag-migrated.}
  \label{fig:tlp-spectrogram-example}
\end{figure}

\subsection{Localised Fourier Shift via Short-Time Fourier Transform}
\label{sec:tlp-stft}

Finally, the three-panel localised-STFT decomposition of
\cref{eq:local-shift,eq:local-phase-line} --- amplitude spectrum
$A(\tau,f)$, phase spectrum $\Dphi(\tau,f)$, and phase-angle domain
$A(\tau,\theta)$ --- is applied with a Gaussian analysis window.
\Cref{fig:tlp-stft-example} shows two of the seven outputs produced by this
sweep; the full set is given in \cref{app:sec-stft}.

\begin{figure}[htbp]
  \centering
  \begin{subfigure}[t]{0.48\textwidth}
    \includegraphics[width=\linewidth]{TLP_078_Localized_Fourier_Shift_STFT_Gaussian_win12_smp____Gazdag.png}
    \caption{Scale 1 of 7}
  \end{subfigure}
  \hfill
  \begin{subfigure}[t]{0.48\textwidth}
    \includegraphics[width=\linewidth]{TLP_081_Localized_Fourier_Shift_STFT_Gaussian_win12_smp____Gazdag.png}
    \caption{Scale 4 of 7}
  \end{subfigure}
  \caption{Localised-STFT phase decomposition (Gaussian window, Gazdag-migrated)
    for two of the seven lateral displacement scales.
    \draftnote{TLP\_078--084 all share an identical auto-generated filename
    stem and do not encode which $\Dx$ each panel corresponds to --- confirm
    the scale for each of the seven figures against the
    \texttt{TimeLapse\_Processing.ipynb} cell order before finalising these
    captions.}}
  \label{fig:tlp-stft-example}
\end{figure}
